\ifdefined\standalone 
\documentclass[12pt]{article}
\usepackage{graphicx}
\usepackage{float}
\usepackage[a4paper, margin=1in]{geometry}
\usepackage[utf8]{inputenc}
\usepackage{textcomp}
\usepackage{amsmath}
\usepackage{booktabs}
\usepackage{tikz}
\usepackage{enumitem}
\usetikzlibrary{decorations.pathmorphing,patterns}
\usepackage[hidelinks]{hyperref} % <- hypertext

\begin{document}
\fi

\section{Thermophysical properties temperature dependancy}

This verification case assesses the ability of \textsc{Gpyro} to correctly evaluate temperature-dependent thermophysical properties. A one-dimensional simulation is performed in which the temperature increases linearly with time at a constant heating rate of $1~\mathrm{K\,s^{-1}}$. The material properties are prescribed as explicit functions of temperature according to the following relations:
\begin{align}
c_p(T)   &= c_{p,0} \left( \frac{T}{T_{\mathrm{ref}}} \right)^{n_c} \\
k(T)   &= k_0 \left( \frac{T}{T_{\mathrm{ref}}} \right)^{n_k}
          + \sigma\,\gamma\,T^3 \\
\rho(T) &= \rho_0 \left( \frac{T}{T_{\mathrm{ref}}} \right)^{n_r}
\end{align}
where $T_{\mathrm{ref}} = 300~\mathrm{K}$ is the reference temperature, $\sigma$ is the Stefan--Boltzmann constant, and the material parameters are listed in Table~\ref{tab:T_Dependant_Props}.
\begin{table}[H]
\centering
\caption{Material parameters used for the temperature-dependent properties verification.}
\label{tab:T_Dependant_Props}
\renewcommand{\arraystretch}{1.3}
\begin{tabular}{c c c c c c c}
\hline
$k_0$ ($\mathrm{W\,m^{-1}\,K^{-1}}$) &
$n_k$ ($-$) &
$\rho_0$ ($\mathrm{kg\,m^{-3}}$) &
$n_r$ ($-$) &
$c_{p,0}$ ($\mathrm{J\,kg^{-1}\,K^{-1}}$) &
$n_c$ ($-$) &
$\gamma$ ($\mathrm{m^{-1}}$) \\
\hline
0.2 & 0.6 & 1430 & 0.25 & 1550 & 0.3 & 0.002 \\
\hline
\end{tabular}
\end{table}

The thermophysical properties predicted by \textsc{Gpyro} are compared against the analytical expressions above as the temperature increases. Results obtained with the original implementation (\textsc{Gpyro} V0.8) are also included, as discrepancies are identified for the original version of \textsc{Gpyro}.

\begin{figure}[H]
\centering
\includegraphics[width=0.8\linewidth]{All_Variables_Subplots_TDep.png}
\caption{Comparison of numerical and analytical temperature-dependent  properties.}
\label{fig:T_evol_fig}
\end{figure}

\ifdefined\standalone
\end{document}
\fi






